\title{2020 年高考数学 全国I卷-文科}
\date{}
\maketitle

\section*{一、选择题}
本题共12小题，每小题5分，共60分。在每小题给出的四个选项中，只有一项是符合题目要求的。

\begin{question}[points=5]
已知集合 $A = \{x \mid x^2 - 3x - 4 < 0\}$, $B = \{-4, 1, 3, 5\}$, 则 $A \cap B =$
\begin{choices}
  \item $\{-4, 1\}$
  \item $\{1, 5\}$
  \item $\{3, 5\}$
  \item $\{1, 3\}$
\end{choices}
\begin{solution}
\textbf{答案：}$\{1, 3\}$

\textbf{分析：}首先解一元二次不等式求得集合 $A$, 之后利用交集中元素的特征求得 $A \cap B$.

\textbf{详解：}由 $x^2 - 3x - 4 < 0$ 解得 $-1 < x < 4$, 所以 $A = \{x \mid -1 < x < 4\}$, 又因为 $B = \{-4, 1, 3, 5\}$, 所以 $A \cap B = \{1, 3\}$.
\end{solution}
\end{question}

\begin{question}[points=5]
若 $z = 1 + 2i + i^3$, 则 $|z| =$
\begin{choices}
  \item $0$
  \item $1$
  \item $\sqrt{2}$
  \item $2$
\end{choices}
\begin{solution}
\textbf{答案：}$\sqrt{2}$

\textbf{分析：}先根据 $i^2 = -1$ 将 $z$ 化简, 再根据向量的模的计算公式即可求出.

\textbf{详解：}因为 $z = 1 + 2i + i^3 = 1 + 2i - i = 1 + i$, 所以 $|z| = \sqrt{1^2 + 1^2} = \sqrt{2}$.
\end{solution}
\end{question}

\begin{question}[points=5]
埃及胡夫金字塔是古代世界建筑奇迹之一，它的形状可视为一个正四棱锥，以该四棱锥的高为边长的正方形面积等于该四棱锥一个侧面三角形的面积，则其侧面三角形底边上的高与底面正方形的边长的比值为
\begin{choices}
  \item $\frac{\sqrt{5} - 1}{4}$
  \item $\frac{\sqrt{5} - 1}{2}$
  \item $\frac{\sqrt{5} + 1}{4}$
  \item $\frac{\sqrt{5} + 1}{2}$
\end{choices}
\begin{solution}
\textbf{答案：}$\frac{\sqrt{5} + 1}{4}$

\textbf{分析：}设 $CD = a$, $PE = b$, 利用 $PO^2 = \frac{1}{2} CD \cdot PE$ 得到关于 $a, b$ 的方程，解方程即可得到答案.

\textbf{详解：}如图，设 $CD = a$, $PE = b$, 则 $PO = \sqrt{PE^2 - OE^2} = \sqrt{b^2 - \frac{a^2}{4}}$. 由题意 $PO^2 = \frac{1}{2} ab$, 即 $b^2 - \frac{a^2}{4} = \frac{1}{2} ab$, 化简得 $4\left(\frac{b}{a}\right)^2 - 2\cdot \frac{b}{a} - 1 = 0$, 解得 $\frac{b}{a} = \frac{1+\sqrt{5}}{4}$ (负值舍去). 故选C.
\end{solution}
\end{question}

\begin{question}[points=5]
设 $O$ 为正方形 $ABCD$ 的中心，在 $O, A, B, C, D$ 中任取 $3$ 点，则取到的 $3$ 点共线的概率为
\begin{choices}
  \item $\frac{1}{5}$
  \item $\frac{2}{5}$
  \item $\frac{1}{2}$
  \item $\frac{4}{5}$
\end{choices}
\begin{solution}
\textbf{答案：}$\frac{1}{5}$

\textbf{分析：}列出从 $5$ 个点选 $3$ 个点的所有情况，再列出 $3$ 点共线的情况，用古典概型的概率计算公式运算即可.

\textbf{详解：}从 $O,A,B,C,D$ 5个点中任取3个有10种不同取法，3点共线只有 $\{A,O,C\}$ 与 $\{B,O,D\}$ 共2种情况，由古典概型的概率计算公式知，取到3点共线的概率为 $\frac{2}{10} = \frac{1}{5}$.
\end{solution}
\end{question}

\begin{question}[points=5]
某校一个课外学习小组为研究某作物种子的发芽率 $y$ 和温度 $x$ (单位：$^\circ$C)的关系，在 $20$ 个不同的温度条件下进行种子发芽实验，由实验数据 $(x_i, y_i)$ ($i=1,2,\cdots,20$) 得到下面的散点图： 由此散点图，在 $10^\circ$C 至 $40^\circ$C 之间，下面四个回归方程类型中最适宜作为发芽率 $y$ 和温度 $x$ 的回归方程类型的是

\begin{center}
  \includegraphics[width=.38\linewidth]{question-5-1.jpg}
\end{center}
\begin{choices}
  \item $y = a + bx$
  \item $y = a + bx^2$
  \item $y = a + be^x$
  \item $y = a + b\ln x$
\end{choices}
\begin{solution}
\textbf{答案：}$y = a + b\ln x$

\textbf{分析：}根据散点图的分布可选择合适的函数模型.

\textbf{详解：}由散点图分布可知，散点图分布在一个对数函数的图象附近，因此，最适合作为发芽率 $y$ 和温度 $x$ 的回归方程类型的是 $y = a + b\ln x$.
\end{solution}
\end{question}

\begin{question}[points=5]
已知圆 $x^2 + y^2 - 6x = 0$, 过点 $(1,2)$ 的直线被该圆所截得的弦的长度的最小值为
\begin{choices}
  \item $1$
  \item $2$
  \item $3$
  \item $4$
\end{choices}
\begin{solution}
\textbf{答案：}$2$

\textbf{分析：}根据直线和圆心与点 $(1,2)$ 连线垂直时，所求的弦长最短，即可得出结论.

\textbf{详解：}圆 $x^2 + y^2 - 6x = 0$ 化为 $(x-3)^2 + y^2 = 9$, 所以圆心 $C$ 坐标为 $(3,0)$, 半径为 $3$. 设 $P(1,2)$, 当过点 $P$ 的直线和直线 $CP$ 垂直时，圆心到过点 $P$ 的直线的距离最大，所求的弦长最短，根据弦长公式最小值为 $2\sqrt{9 - |CP|^2} = 2\sqrt{9 - 8} = 2$.
\end{solution}
\end{question}

\begin{question}[points=5]
设函数 $f(x) = \cos\left(\omega x + \frac{\pi}{6}\right)$ 在 $[-\pi,\pi]$ 的图像大致如下图，则 $f(x)$ 的最小正周期为

\begin{center}
  \includegraphics[width=.38\linewidth]{question-7-1.jpg}
\end{center}
\begin{choices}
  \item $\frac{10\pi}{9}$
  \item $\frac{7\pi}{6}$
  \item $\frac{4\pi}{3}$
  \item $\frac{3\pi}{2}$
\end{choices}
\begin{solution}
\textbf{答案：}$\frac{4\pi}{3}$

\textbf{分析：}由图可得函数图象过点 $\left(-\frac{4\pi}{9},0\right)$, 代入得到 $\cos\left(-\frac{4\pi}{9}\omega + \frac{\pi}{6}\right) = 0$, 结合点 $\left(-\frac{4\pi}{9},0\right)$ 是函数图象与 $x$ 轴负半轴的第一个交点得到 $-\frac{4\pi}{9}\omega + \frac{\pi}{6} = -\frac{\pi}{2}$, 解得 $\omega = \frac{3}{2}$, 再利用周期公式即可得解.

\textbf{详解：}由图可得函数图象过点 $\left(-\frac{4\pi}{9},0\right)$, 代入得 $\cos\left(-\frac{4\pi}{9}\omega + \frac{\pi}{6}\right) = 0$. 点 $\left(-\frac{4\pi}{9},0\right)$ 是函数图象与 $x$ 轴负半轴的第一个交点，所以 $-\frac{4\pi}{9}\omega + \frac{\pi}{6} = -\frac{\pi}{2}$, 解得 $\omega = \frac{3}{2}$. 所以函数 $f(x)$ 的最小正周期为 $T = \frac{2\pi}{\omega} = \frac{2\pi}{3/2} = \frac{4\pi}{3}$.
\end{solution}
\end{question}

\begin{question}[points=5]
设 $a\log_3 4 = 2$, 则 $4^{-a} =$
\begin{choices}
  \item $\frac{1}{16}$
  \item $\frac{1}{9}$
  \item $\frac{1}{8}$
  \item $\frac{1}{6}$
\end{choices}
\begin{solution}
\textbf{答案：}$\frac{1}{9}$

\textbf{分析：}首先根据题中所给的式子，结合对数的运算法则，得到 $\log_3 4^a = 2$, 即 $4^a = 9$, 进而求得 $4^{-a} = \frac{1}{9}$.

\textbf{详解：}由 $a\log_3 4 = 2$ 可得 $\log_3 4^a = 2$, 所以 $4^a = 9$, 所以 $4^{-a} = \frac{1}{9}$.
\end{solution}
\end{question}

\begin{question}[points=5]
执行下面的程序框图，则输出的 $n=$

\begin{center}
  \includegraphics[width=.38\linewidth]{question-9-1.jpg}
\end{center}
\begin{choices}
  \item $17$
  \item $19$
  \item $21$
  \item $23$
\end{choices}
\begin{solution}
\textbf{答案：}$21$

\textbf{分析：}根据程序框图的算法功能可知，要计算满足 $1+3+5+\cdots+n > 100$ 的最小正奇数 $n$, 根据等差数列求和公式即可求出.

\textbf{详解：}输出的 $n$ 是满足 $1+3+5+\cdots+n > 100$ 的最小正奇数，因为 $1+3+5+\cdots+n = \frac{(1+n)\times(\frac{n-1}{2}+1)}{2} = \frac{(n+1)^2}{4} > 100$, 解得 $n > 19$, 所以输出的 $n=21$.
\end{solution}
\end{question}

\begin{question}[points=5]
设 $\{a_n\}$ 是等比数列，且 $a_1 + a_2 + a_3 = 1$, $a_2 + a_3 + a_4 = 2$, 则 $a_6 + a_7 + a_8 =$
\begin{choices}
  \item $12$
  \item $24$
  \item $30$
  \item $32$
\end{choices}
\begin{solution}
\textbf{答案：}$32$

\textbf{分析：}根据已知条件求得 $q$ 的值，再由 $a_6+a_7+a_8 = q^5(a_1+a_2+a_3)$ 可求得结果.

\textbf{详解：}设等比数列 $\{a_n\}$ 的公比为 $q$, 则 $a_1+a_2+a_3 = a_1(1+q+q^2)=1$, $a_2+a_3+a_4 = a_1q(1+q+q^2)=q=2$, 因此 $a_6+a_7+a_8 = a_1q^5(1+q+q^2)=q^5=32$.
\end{solution}
\end{question}

\begin{question}[points=5]
设 $F_1, F_2$ 是双曲线 $C: x^2 - \frac{y^2}{3} = 1$ 的两个焦点，$O$ 为坐标原点，点 $P$ 在 $C$ 上且 $|OP| = 2$, 则 $\triangle PF_1F_2$ 的面积为
\begin{choices}
  \item $\frac{7}{2}$
  \item $3$
  \item $\frac{5}{2}$
  \item $2$
\end{choices}
\begin{solution}
\textbf{答案：}$3$

\textbf{分析：}由 $|OP| = 2$ 得点 $P$ 在以 $F_1F_2$ 为直径的圆上，即 $\triangle PF_1F_2$ 是以 $P$ 为直角顶点的直角三角形，利用勾股定理和双曲线定义求出 $|PF_1||PF_2|$, 代入面积公式计算.

\textbf{详解：}由已知，$F_1(-2,0), F_2(2,0)$, 则 $a=1, c=2$, 因为 $|OP|=2 = \frac{1}{2}|F_1F_2|$, 所以点 $P$ 在以 $F_1F_2$ 为直径的圆上，即 $\triangle PF_1F_2$ 是以 $P$ 为直角顶点的直角三角形，故 $|PF_1|^2 + |PF_2|^2 = |F_1F_2|^2 = 16$, 又 $|PF_1|-|PF_2|=2a=2$, 所以 $4 = (|PF_1|-|PF_2|)^2 = |PF_1|^2 + |PF_2|^2 - 2|PF_1||PF_2| = 16 - 2|PF_1||PF_2|$, 解得 $|PF_1||PF_2|=6$, 所以 $S_{\triangle PF_1F_2} = \frac{1}{2}|PF_1||PF_2|=3$.
\end{solution}
\end{question}

\begin{question}[points=5]
已知 $A, B, C$ 为球 $O$ 的球面上的三个点，$\odot O_1$ 为 $\triangle ABC$ 的外接圆，若 $\odot O_1$ 的面积为 $4\pi$, $AB = BC = AC = OO_1$, 则球 $O$ 的表面积为
\begin{choices}
  \item $64\pi$
  \item $48\pi$
  \item $36\pi$
  \item $32\pi$
\end{choices}
\begin{solution}
\textbf{答案：}$64\pi$

\textbf{分析：}由已知可得等边 $\triangle ABC$ 的外接圆半径，进而求出其边长，得出 $OO_1$ 的值，根据球截面性质，求出球的半径，即可得出结论.

\textbf{详解：}设圆 $O_1$ 半径为 $r$, 球的半径为 $R$, 依题意得 $\pi r^2 = 4\pi$, 所以 $r=2$. 由正弦定理可得 $AB = 2r\sin 60^\circ = 2\sqrt{3}$, 所以 $OO_1 = AB = 2\sqrt{3}$. 根据圆截面性质 $OO_1 \perp$ 平面 $ABC$, 所以 $OO_1 \perp O_1A$, $R = OA = \sqrt{OO_1^2 + O_1A^2} = \sqrt{(2\sqrt{3})^2 + 2^2} = 4$, 所以球 $O$ 的表面积 $S = 4\pi R^2 = 64\pi$.
\end{solution}
\end{question}

\section*{二、填空题}
本题共4小题，每小题5分，共20分。

\begin{question}[points=5]
若 $x, y$ 满足约束条件 $\begin{cases} 2x + y - 2 \le 0, \\ x - y - 1 \ge 0, \\ y + 1 \ge 0, \end{cases}$ 则 $z = x + 7y$ 的最大值为
\fillin[{$1$}]
\begin{solution}
\textbf{答案：}$1$

\textbf{分析：}首先画出可行域，然后结合目标函数的几何意义即可求得其最大值.

\textbf{详解：}绘制不等式组表示的平面区域如图所示，目标函数 $z=x+7y$ 即 $y = -\frac{1}{7}x + \frac{1}{7}z$, 其中 $z$ 取得最大值时表示直线系在 $y$ 轴上的截距最大，据此结合目标函数的几何意义可知目标函数在点 $A$ 处取得最大值，联立直线方程 $\begin{cases} 2x+y-2=0 \\ x-y-1=0 \end{cases}$ 可得点 $A(1,0)$, 所以 $z_{\max}=1+7\times0=1$.
\end{solution}
\end{question}

\begin{question}[points=5]
设向量 $\vec{a} = (1, -1)$, $\vec{b} = (m+1, 2m-4)$, 若 $\vec{a} \perp \vec{b}$, 则 $m =$
\fillin[{$5$}]
\begin{solution}
\textbf{答案：}$5$

\textbf{分析：}根据向量垂直的坐标表示求得结果.

\textbf{详解：}由 $\vec{a} \perp \vec{b}$ 可得 $\vec{a} \cdot \vec{b} = 0$, 即 $1\cdot(m+1) + (-1)\cdot(2m-4) = 0$, 解得 $m=5$.
\end{solution}
\end{question}

\begin{question}[points=5]
曲线 $y = \ln x + x + 1$ 的一条切线的斜率为 $2$, 则该切线的方程为
\fillin[{$y = 2x$}]
\begin{solution}
\textbf{答案：}$y = 2x$

\textbf{分析：}设切线的切点坐标为 $(x_0, y_0)$, 对函数求导，利用 $y'|_{x_0} = 2$, 求出 $x_0$, 代入曲线方程求出 $y_0$, 得到切线的点斜式方程，化简即可.

\textbf{详解：}设切点坐标为 $(x_0, y_0)$, $y' = \frac{1}{x} + 1$, 则 $y'|_{x=x_0} = \frac{1}{x_0} + 1 = 2$, 解得 $x_0 = 1$, 代入得 $y_0 = \ln 1 + 1 + 1 = 2$, 所以切线方程为 $y - 2 = 2(x - 1)$, 即 $y = 2x$.
\end{solution}
\end{question}

\begin{question}[points=5]
数列 $\{a_n\}$ 满足 $a_{n+2} + (-1)^n a_n = 3n - 1$, 前 $16$ 项和为 $540$, 则 $a_1 =$
\fillin[{$7$}]
\begin{solution}
\textbf{答案：}$7$

\textbf{分析：}对 $n$ 为奇偶数分类讨论，分别得出奇数项、偶数项的递推关系，由奇数项递推公式将奇数项用 $a_1$ 表示，由偶数项递推公式得出偶数项的和，建立 $a_1$ 方程，求解即可得出结论.

\textbf{详解：}由 $a_{n+2} + (-1)^n a_n = 3n - 1$, 当 $n$ 为奇数时，$a_{n+2} = a_n + 3n - 1$; 当 $n$ 为偶数时，$a_{n+2} + a_n = 3n - 1$. 设数列 $\{a_n\}$ 的前 $n$ 项和为 $S_n$, $S_{16} = a_1 + a_2 + \cdots + a_{16} = a_1 + a_3 + a_5 + \cdots + a_{15} + (a_2 + a_4) + \cdots + (a_{14} + a_{16}) = a_1 + (a_1+2) + (a_1+10) + (a_1+24) + (a_1+44) + (a_1+70) + (a_1+102) + (a_1+140) + (5+17+29+41) = 8a_1 + 392 + 92 = 8a_1 + 484 = 540$, 解得 $a_1 = 7$.
\end{solution}
\end{question}

\section*{三、解答题}
共70分。解答应写出文字说明、证明过程或演算步骤。第17~21题为必考题，每个试题考生都必须作答。第22、23题为选考题，考生根据要求作答。

\begin{question}[points=12]
某厂接受了一项加工业务，加工出来的产品(单位：件)按标准分为 $A, B, C, D$ 四个等级。加工业务约定：对于 $A$ 级品、$B$ 级品、$C$ 级品，厂家每件分别收取加工费 $90$ 元，$50$ 元，$20$ 元；对于 $D$ 级品，厂家每件要赔偿原料损失费 $50$ 元。该厂有甲、乙两个分厂可承接加工业务。甲分厂加工成本费为 $25$ 元/件，乙分厂加工成本费为 $20$ 元/件。厂家为决定由哪个分厂承接加工业务，在两个分厂各试加工了 $100$ 件这种产品，并统计了这些产品的等级，整理如下：



甲分厂产品等级的频数分布表



| 等级 | A | B | C | D |

|------|---|---|---|---|

| 频数 | 40 | 20 | 20 | 20 |



乙分厂产品等级的频数分布表



| 等级 | A | B | C | D |

|------|---|---|---|---|

| 频数 | 28 | 17 | 34 | 21 |



(1) 分别估计甲、乙两分厂加工出来的一件产品为 $A$ 级品的概率；



(2) 分别求甲、乙两分厂加工出来的 $100$ 件产品的平均利润，以平均利润为依据，厂家应选哪个分厂承接加工业务?
\begin{solution}
\textbf{答案：}(1) 甲分厂加工出来的 $A$ 级品的概率为 $0.4$, 乙分厂加工出来的 $A$ 级品的概率为 $0.28$; (2) 选甲分厂.

\textbf{分析：}(1) 根据两个频数分布表即可求出； (2) 根据题意分别求出甲乙两厂加工 $100$ 件产品的总利润，即可求出平均利润，由此作出选择.

\textbf{详解：}(1) 由表可知，甲厂加工出来的一件产品为 $A$ 级品的概率为 $\frac{40}{100}=0.4$, 乙厂加工出来的一件产品为 $A$ 级品的概率为 $\frac{28}{100}=0.28$.

(2) 甲分厂加工 $100$ 件产品的总利润为 $40\times(90-25)+20\times(50-25)+20\times(20-25)-20\times(50+25)=1500$ 元, 平均利润为 $15$ 元每件; 乙分厂加工 $100$ 件产品的总利润为 $28\times(90-20)+17\times(50-20)+34\times(20-20)-21\times(50+20)=1000$ 元, 平均利润为 $10$ 元每件. 故厂家选择甲分厂承接加工任务.
\end{solution}
\end{question}

\begin{question}[points=12]
$\triangle ABC$ 的内角 $A, B, C$ 的对边分别为 $a, b, c$. 已知 $B=150^\circ$.



(1) 若 $a=\sqrt{3}c$, $b=2\sqrt{7}$, 求 $\triangle ABC$ 的面积;



(2) 若 $\sin A + \sqrt{3}\sin C = \frac{\sqrt{2}}{2}$, 求 $C$.
\begin{solution}
\textbf{答案：}(1) $\sqrt{3}$; (2) $15^\circ$.

\textbf{分析：}(1) 已知角 $B$ 和 $b$ 边，结合 $a,c$ 关系，由余弦定理建立 $c$ 的方程，求解得出 $a,c$, 利用面积公式，即可得出结论； (2) 将 $A=30^\circ - C$ 代入已知等式，由两角差的正弦和辅助角公式，化简得出有关 $C$ 角的三角函数值，结合 $C$ 的范围，即可求解.

\textbf{详解：}(1) 由余弦定理可得 $b^2 = 28 = a^2 + c^2 - 2ac\cos150^\circ = 7c^2$, 所以 $c=2$, $a=2\sqrt{3}$, $\triangle ABC$ 的面积 $S = \frac{1}{2}ac\sin B = \sqrt{3}$.

(2) 因为 $A+C=30^\circ$, 所以 $\sin A + \sqrt{3}\sin C = \sin(30^\circ - C) + \sqrt{3}\sin C = \frac{1}{2}\cos C + \frac{\sqrt{3}}{2}\sin C = \sin(C+30^\circ) = \frac{\sqrt{2}}{2}$. 因为 $0^\circ < C < 30^\circ$, 所以 $30^\circ < C+30^\circ < 60^\circ$, 所以 $C+30^\circ = 45^\circ$, 即 $C=15^\circ$.
\end{solution}
\end{question}

\begin{question}[points=12]
如图，$D$ 为圆锥的顶点，$O$ 是圆锥底面的圆心，$\triangle ABC$ 是底面的内接正三角形，$P$ 为 $DO$ 上一点，$\angle APC = 90^\circ$.



(1) 证明：平面 $PAB \perp$ 平面 $PAC$;



(2) 设 $DO = \sqrt{2}$, 圆锥的侧面积为 $\sqrt{3}\pi$, 求三棱锥 $P-ABC$ 的体积.

\begin{center}
  \includegraphics[width=.38\linewidth]{question-19-1.jpg}
\end{center}
\begin{solution}
\textbf{答案：}(1) 证明见解析; (2) $\frac{\sqrt{6}}{8}$.

\textbf{分析：}(1) 根据已知可得 $PA=PB=PC$, 进而有 $\triangle PAC \cong \triangle PBC$, 可得 $\angle APC = \angle BPC = 90^\circ$, 即 $PB \perp PC$, 从而证得 $PC \perp$ 平面 $PAB$, 即可证得结论； (2) 将已知条件转化为母线 $l$ 和底面半径 $r$ 的关系，进而求出底面半径，由正弦定理，求出正三角形 $ABC$ 边长，在等腰直角三角形 $APC$ 中求出 $AP$, 在 $\triangle APO$ 中求出 $PO$, 即可求出结论.

\textbf{详解：}(1) 因为 $D$ 为圆锥顶点，$O$ 为底面圆心，所以 $OD \perp$ 平面 $ABC$. 因为 $P$ 在 $DO$ 上，$OA = OB = OC$, 所以 $PA = PB = PC$. 又 $\triangle ABC$ 是圆内接正三角形，所以 $AC = BC$, 故 $\triangle PAC \cong \triangle PBC$, 所以 $\angle APC = \angle BPC = 90^\circ$, 即 $PB \perp PC$, $PA \perp PC$. 因为 $PA \cap PB = P$, 所以 $PC \perp$ 平面 $PAB$, 又 $PC \subset$ 平面 $PAC$, 所以平面 $PAB \perp$ 平面 $PAC$.

(2) 设圆锥的母线为 $l$, 底面半径为 $r$, 圆锥的侧面积为 $\pi rl = \sqrt{3}\pi$, 所以 $rl = \sqrt{3}$. 又 $OD^2 = l^2 - r^2 = 2$, 解得 $r=1$, $l=\sqrt{3}$. 正三角形 $ABC$ 的边长 $AC = 2r\sin60^\circ = \sqrt{3}$. 在等腰直角三角形 $APC$ 中，$AP = \frac{\sqrt{2}}{2} AC = \frac{\sqrt{6}}{2}$. 在 $\triangle APO$ 中，$PO = \sqrt{AP^2 - OA^2} = \sqrt{\frac{6}{4} - 1} = \frac{\sqrt{2}}{2}$. 所以三棱锥 $P-ABC$ 的体积 $V_{P-ABC} = \frac{1}{3} PO \cdot S_{\triangle ABC} = \frac{1}{3} \times \frac{\sqrt{2}}{2} \times \frac{\sqrt{3}}{4} \times \sqrt{3} = \frac{\sqrt{6}}{8}$.
\end{solution}
\end{question}

\begin{question}[points=12]
已知函数 $f(x) = e^x - a(x+2)$.



(1) 当 $a=1$ 时，讨论 $f(x)$ 的单调性;



(2) 若 $f(x)$ 有两个零点，求 $a$ 的取值范围.
\begin{solution}
\textbf{答案：}(1) 减区间为 $(-\infty, 0)$, 增区间为 $(0, +\infty)$; (2) $\left(\frac{1}{e}, +\infty\right)$.

\textbf{分析：}(1) 将 $a=1$ 代入函数解析式，对函数求导，分别令导数大于零和小于零，求得函数的单调增区间和减区间； (2) 若 $f(x)$ 有两个零点，即 $e^x - a(x+2)=0$ 有两个解，将其转化为 $a = \frac{e^x}{x+2}$ 有两个解，令 $h(x)=\frac{e^x}{x+2} (x\neq -2)$, 求导研究函数图象的走向，从而求得结果.

\textbf{详解：}(1) 当 $a=1$ 时，$f(x)=e^x-(x+2)$, $f'(x)=e^x-1$, 令 $f'(x)<0$ 解得 $x<0$, 令 $f'(x)>0$ 解得 $x>0$, 所以 $f(x)$ 的减区间为 $(-\infty,0)$, 增区间为 $(0,+\infty)$.

(2) 若 $f(x)$ 有两个零点，即 $e^x - a(x+2)=0$ 有两个解，从方程可知 $x=-2$ 不成立，即 $a = \frac{e^x}{x+2}$ 有两个解. 令 $h(x)=\frac{e^x}{x+2} (x\neq -2)$, 则 $h'(x)=\frac{e^x(x+2)-e^x}{(x+2)^2}=\frac{e^x(x+1)}{(x+2)^2}$. 令 $h'(x)>0$ 解得 $x>-1$, 令 $h'(x)<0$ 解得 $x<-2$ 或 $-2<x<-1$, 所以 $h(x)$ 在 $(-\infty,-2)$ 和 $(-2,-1)$ 上单调递减，在 $(-1,+\infty)$ 上单调递增，且当 $x<-2$ 时 $h(x)<0$, 当 $x\to -2^+$ 时 $h(x)\to +\infty$, 当 $x\to +\infty$ 时 $h(x)\to +\infty$. 所以当 $a = \frac{e^x}{x+2}$ 有两个解时，有 $a > h(-1) = \frac{1}{e}$, 所以 $a$ 的取值范围是 $\left(\frac{1}{e}, +\infty\right)$.
\end{solution}
\end{question}

\begin{question}[points=12]
已知 $A, B$ 分别为椭圆 $E: \frac{x^2}{a^2} + y^2 = 1$ ($a>1$) 的左、右顶点，$G$ 为 $E$ 的上顶点，$\overrightarrow{AG} \cdot \overrightarrow{GB} = 8$, $P$ 为直线 $x=6$ 上的动点，$PA$ 与 $E$ 的另一交点为 $C$, $PB$ 与 $E$ 的另一交点为 $D$.



(1) 求 $E$ 的方程;



(2) 证明：直线 $CD$ 过定点.
\begin{solution}
\textbf{答案：}(1) $\frac{x^2}{9} + y^2 = 1$; (2) 证明见解析, 定点为 $\left(\frac{3}{2}, 0\right)$.

\textbf{分析：}(1) 由已知可得 $A(-a,0)$, $B(a,0)$, $G(0,1)$, 即可求得 $\overrightarrow{AG}\cdot\overrightarrow{GB}=a^2-1$, 结合已知即可求得 $a^2=9$; (2) 设 $P(6,y_0)$, 联立直线 $AP$ 与椭圆方程求得点 $C$ 坐标，同理求得点 $D$ 坐标，表示出直线 $CD$ 的方程，整理可得直线 $CD$ 过定点 $\left(\frac{3}{2},0\right)$.

\textbf{详解：}(1) 由椭圆方程得 $A(-a,0)$, $B(a,0)$, $G(0,1)$, 则 $\overrightarrow{AG}=(a,1)$, $\overrightarrow{GB}=(a,-1)$, 所以 $\overrightarrow{AG}\cdot\overrightarrow{GB}=a^2-1=8$, 解得 $a^2=9$, 故椭圆方程为 $\frac{x^2}{9}+y^2=1$.

(2) 设 $P(6,y_0)$, 则直线 $AP$ 的方程为 $y=\frac{y_0}{9}(x+3)$. 联立 $\begin{cases} \frac{x^2}{9}+y^2=1 \\ y=\frac{y_0}{9}(x+3) \end{cases}$, 消 $y$ 得 $(y_0^2+9)x^2+6y_0^2x+9y_0^2-81=0$, 由韦达定理知 $x_A x_C = \frac{9y_0^2-81}{y_0^2+9}$, 又 $x_A=-3$, 所以 $x_C = \frac{-3y_0^2+27}{y_0^2+9}$, 代入直线 $AP$ 得 $y_C = \frac{6y_0}{y_0^2+9}$, 即 $C\left(\frac{-3y_0^2+27}{y_0^2+9}, \frac{6y_0}{y_0^2+9}\right)$. 同理，直线 $PB$ 的方程为 $y=\frac{y_0}{3}(x-3)$, 联立解得 $D\left(\frac{3y_0^2-3}{y_0^2+1}, \frac{-2y_0}{y_0^2+1}\right)$. 则直线 $CD$ 的斜率 $k_{CD} = \frac{\frac{6y_0}{y_0^2+9} - \frac{-2y_0}{y_0^2+1}}{\frac{-3y_0^2+27}{y_0^2+9} - \frac{3y_0^2-3}{y_0^2+1}} = \frac{8y_0(y_0^2+3)}{6(9-y_0^4)} = \frac{4y_0}{3(3-y_0^2)}$, 所以直线 $CD$ 的方程为 $y - \frac{-2y_0}{y_0^2+1} = \frac{4y_0}{3(3-y_0^2)}\left(x - \frac{3y_0^2-3}{y_0^2+1}\right)$, 化简得 $y = \frac{4y_0}{3(3-y_0^2)}\left(x - \frac{3}{2}\right)$, 所以直线 $CD$ 过定点 $\left(\frac{3}{2},0\right)$.
\end{solution}
\end{question}

\begin{question}[points=10]
在直角坐标系 $xOy$ 中，曲线 $C_1$ 的参数方程为 $\begin{cases} x = \cos^k t, \\ y = \sin^k t \end{cases}$ ($t$ 为参数). 以坐标原点为极点，$x$ 轴正半轴为极轴建立极坐标系，曲线 $C_2$ 的极坐标方程为 $4\rho\cos\theta - 16\rho\sin\theta + 3 = 0$.



(1) 当 $k=1$ 时，$C_1$ 是什么曲线?



(2) 当 $k=4$ 时，求 $C_1$ 与 $C_2$ 的公共点的直角坐标.
\begin{solution}
\textbf{答案：}(1) 曲线 $C_1$ 表示以坐标原点为圆心，半径为 $1$ 的圆； (2) $\left(\frac{1}{4}, \frac{1}{4}\right)$.

\textbf{分析：}(1) 利用 $\sin^2 t + \cos^2 t = 1$ 消去参数 $t$, 求出曲线 $C_1$ 的普通方程，即可得出结论； (2) 当 $k=4$ 时，$x\ge 0, y\ge 0$, 曲线 $C_1$ 的参数方程化为 $\begin{cases} x = \cos^2 t, \\ y = \sin^2 t \end{cases}$, 两式相加消去参数 $t$, 得 $C_1$ 普通方程，由 $\rho\cos\theta = x, \rho\sin\theta = y$ 将曲线 $C_2$ 化为直角坐标方程，联立 $C_1, C_2$ 方程，即可求解.

\textbf{详解：}(1) 当 $k=1$ 时，曲线 $C_1$ 的参数方程为 $\begin{cases} x = \cos t, \\ y = \sin t \end{cases}$ ($t$ 为参数), 两式平方相加得 $x^2+y^2=1$, 所以曲线 $C_1$ 表示以坐标原点为圆心，半径为 $1$ 的圆.

(2) 当 $k=4$ 时，曲线 $C_1$ 的参数方程为 $\begin{cases} x = \cos^4 t, \\ y = \sin^4 t \end{cases}$ ($t$ 为参数), 所以 $x\ge0, y\ge0$, 化为 $\begin{cases} x = \cos^2 t, \\ y = \sin^2 t \end{cases}$, 两式相加得 $\sqrt{x} + \sqrt{y} = 1$, 平方得 $y = x - 2\sqrt{x} + 1$, $0\le x\le 1$, $0\le y\le 1$. 曲线 $C_2$ 的直角坐标方程为 $4x - 16y + 3 = 0$. 联立 $\begin{cases} y = x - 2\sqrt{x} + 1 \\ 4x - 16y + 3 = 0 \end{cases}$, 消 $y$ 得 $12x - 32\sqrt{x} + 13 = 0$, 解得 $\sqrt{x} = \frac{1}{2}$ 或 $\sqrt{x} = \frac{13}{6}$ (舍去), 所以 $x = \frac{1}{4}$, $y = \frac{1}{4}$, 故公共点的直角坐标为 $\left(\frac{1}{4}, \frac{1}{4}\right)$.
\end{solution}
\end{question}

\begin{question}[points=10]
已知函数 $f(x) = |3x+1| - 2|x-1|$.



(1) 画出 $y = f(x)$ 的图像;



(2) 求不等式 $f(x) > f(x+1)$ 的解集.

\begin{center}
  \includegraphics[width=.38\linewidth]{question-23-1.jpg}
\end{center}
\begin{solution}
\textbf{答案：}(1) 图像见解析； (2) $\left(-\infty, -\frac{7}{6}\right)$.

\textbf{分析：}(1) 根据分段讨论法，即可写出函数 $f(x)$ 的解析式，作出图象； (2) 作出函数 $f(x+1)$ 的图象，根据图象即可解出.

\textbf{详解：}(1) $f(x) = \begin{cases} x+3, & x\ge 1, \\ 5x-1, & -\frac{1}{3} < x < 1, \\ -x-3, & x\le -\frac{1}{3}. \end{cases}$ 图象如图所示.

(2) 将函数 $f(x)$ 的图象向左平移 $1$ 个单位，可得函数 $f(x+1)$ 的图象，如图所示. 由 $-x-3 = 5(x+1)-1$ 解得 $x = -\frac{7}{6}$. 所以不等式的解集为 $\left(-\infty, -\frac{7}{6}\right)$.
\end{solution}
\end{question}

